% Generated draft: requires mathematical and source review.
\documentclass[11pt]{article}
\usepackage[margin=25mm]{geometry}
\usepackage{fontspec}
\setmainfont{Latin Modern Roman}
\usepackage{amsmath,amssymb,mathtools,unicode-math}
\setmathfont{Latin Modern Math}
\usepackage{enumitem,needspace,xurl,hyperref}
\setlength{\parindent}{0pt}
\setlength{\parskip}{5pt}
\setlength{\emergencystretch}{3em}
\newcommand{\N}{\mathbb{N}}
\newcommand{\Z}{\mathbb{Z}}
\newcommand{\Q}{\mathbb{Q}}
\newcommand{\R}{\mathbb{R}}
\newcommand{\C}{\mathbb{C}}
\newcommand{\F}{\mathbb{F}}
\newcommand{\A}{\mathbb{A}}
\newcommand{\PP}{\mathbb{P}}
\newcommand{\E}{\mathbb{E}}
\newcommand{\eps}{\varepsilon}
\newcommand{\dd}{\,\mathrm{d}}
\begin{document}
\section*{UnsolvedMath ID: 10400045}
\textbf{Source problem code:} catalogue entry
\subsection{Definitions and mathematical statement}
Let \(\mathcal K\) denote the source's knot class, let \(O\) be the unknot, and let \(\sim_{C_d}\) be its \(C_d\)-equivalence relation. Under the usual interpretation these are knots in \(S^3\) and equivalences generated by the specified degree-\(d\) clasper moves, but the record provides no move definition or knot-category conventions. For each admissible index \(d\), form
\[
 G_d=\{K\in\mathcal K:K\sim_{C_d}O\}/\sim_{C_{d+1}}.
\]
Conjecture 2.23 states that \(G_d\) is torsion free for every \(d\). With the source's intended group operation of connected sum and identity \([O]\), this means
\[
 \forall [K]\in G_d\ \forall m\in\Z_{>0}:\quad
 m[K]=[O]\ \Longrightarrow\ [K]=[O].
\]
Equivalently, under that group interpretation, for every \(K\sim_{C_d}O\) and every positive integer \(m\),
\[
 K^{\# m}\sim_{C_{d+1}}O\ \Longrightarrow\ K\sim_{C_{d+1}}O,
\]
where \(K^{\#m}\) is the connected sum of \(m\) copies of \(K\). The definition and nesting of the \(C_d\)-relations, the induced group structure, and the index range must be supplied from the source; ``for each \(d\)'' is retained without inventing a \(d=0\) convention.
\subsection{Short English statement}
At each level of the \(C_d\)-equivalence filtration, the group of knots equivalent to the unknot at that level, modulo the next equivalence, is conjectured to have no nonzero torsion. Thus a finite connected sum becoming trivial at the next level should force the original class to be trivial there as well.
\subsection{Sources}
\begin{itemize}
\item Local UnsolvedMath registry, permanent ID \(10400045\), exact source code \texttt{catalogue entry}, Conjecture 2.23. No individual proposer or proposed year is supplied.
\item Ohtsuki, \emph{Problems on invariants of knots and 3-manifolds} (2002), Conjecture 2.23, PDF page 47: \url{https://msp.org/gtm/2002/04/gtm-2002-04-024s.pdf}. Printed bold headings distinguish 231 numbered items from references to them in exposition.
\item Extraction pdf-font-structure, accessed 2026-07-29. Source evidence NEEDS\_REVIEW presents an open problem and requires a current primary-literature audit. The supplied OPEN-TRIAGE summary calls this the torsion-free conjecture for \(C_d\)-equivalence and says it is not obviously resolved; it supplies no further citation.
\item Dataset attribution: Ulam AI / UnsolvedMath Contributors, CC BY 4.0. Source rights NEEDS\_REVIEW: publicly accessible author/publisher copy, redistribution terms requiring review.
\end{itemize}
\end{document}
